% نقش: صفحه‌آرایی، عنوان‌ها، فهرست‌ها، قضیه‌ها، کد و نمودار.
% محتوای علمی پروژه را اینجا ننویسید.

% متن ۱۴ پوینت با فاصلهٔ خط ۱٫۵ (پایه ۲۱ پوینت) مطابق دستورالعمل دانشگاه
\AtBeginDocument{\fontsize{14pt}{21pt}\selectfont}
\setlength{\parindent}{1.25cm}
\setlength{\parskip}{0pt}
\setlength{\emergencystretch}{2em}
\raggedbottom

% شماره صفحه وسط پایین؛ fancyhdr با bidi تداخل می‌سازد.
\pagestyle{plain}

\titleformat{\chapter}[display]
  {\centering\TitleFont\bfseries\fontsize{16pt}{24pt}\selectfont}
  {فصل \thechapter}
  {0.5em}
  {}
\titleformat{name=\chapter,numberless}[display]
  {\centering\TitleFont\bfseries\fontsize{16pt}{24pt}\selectfont}
  {}
  {0pt}
  {}
\titlespacing*{\chapter}{0pt}{-10pt}{28pt}

\titleformat{\section}
  {\HeadingFont\bfseries\itshape\fontsize{14pt}{21pt}\selectfont}
  {\thesection}
  {0.75em}
  {}
\titlespacing*{\section}{0pt}{21pt}{7pt}

\titleformat{\subsection}
  {\HeadingFont\bfseries\itshape\fontsize{14pt}{21pt}\selectfont}
  {\thesubsection}
  {0.75em}
  {}
\titlespacing*{\subsection}{0pt}{14pt}{5pt}

\titleformat{\subsubsection}
  {\HeadingFont\bfseries\fontsize{13pt}{19pt}\selectfont}
  {\thesubsubsection}
  {0.75em}
  {}

% جداکنندهٔ شمارهٔ فصل‌محور مطابق دستورالعمل دانشگاه
\SepMark{-}

% رابطه از چپ و شماره از راست (کلاس fleqn)
\setlength{\mathindent}{0pt}

\DeclareCaptionFont{bonabcaption}{\fontsize{10pt}{15pt}\selectfont}
\captionsetup{
  font={bonabcaption,stretch=1.15},
  labelfont=bf,
  justification=centering,
  singlelinecheck=false,
  labelsep=period
}
\captionsetup[table]{position=top,skip=7pt}
\captionsetup[figure]{position=bottom,skip=7pt}

% گلوله از قلم ریاضی؛ B Mitra نویسهٔ bullet ندارد.
\setlist[itemize]{label={\ensuremath{\bullet}},topsep=4pt,itemsep=2pt,parsep=0pt}
\setlist[enumerate]{topsep=4pt,itemsep=2pt,parsep=0pt}

\theoremstyle{definition}
\newtheorem{definition}{تعریف}[chapter]
\newtheorem{example}[definition]{مثال}
\theoremstyle{plain}
\newtheorem{theorem}[definition]{قضیه}
\newtheorem{lemma}[definition]{لم}
\newtheorem{proposition}[definition]{گزاره}
\theoremstyle{remark}
\newtheorem{remark}[definition]{ملاحظه}
\renewcommand{\proofname}{برهان}

\definecolor{codebg}{RGB}{247,248,250}
\definecolor{codecomment}{RGB}{45,125,70}
\definecolor{codekeyword}{RGB}{25,65,150}
\definecolor{codestring}{RGB}{160,65,45}
\renewcommand{\lstlistingname}{{\HeadingFont کد}}
\renewcommand{\lstlistlistingname}{فهرست کدها}
\lstdefinestyle{bonabcode}{
  basicstyle=\ttfamily\footnotesize,
  keywordstyle=\color{codekeyword}\bfseries,
  commentstyle=\color{codecomment},
  stringstyle=\color{codestring},
  backgroundcolor=\color{codebg},
  frame=single,
  rulecolor=\color{black!20},
  numbers=left,
  numberstyle=\tiny\color{black!55},
  numbersep=8pt,
  showstringspaces=false,
  breaklines=true,
  tabsize=4,
  captionpos=b,
  xleftmargin=1.5em,
  framexleftmargin=1.2em
}
\lstdefinestyle{bonabpython}{style=bonabcode,language=Python}
\lstset{style=bonabpython}

% سبک‌های کمکی نمودار؛ در فصل‌ها با نام tikzstyle استفاده شوند.
\tikzset{
  block/.style={draw,rounded corners,minimum width=3.2cm,minimum height=1.05cm,align=center,fill=blue!5},
  loss/.style={draw,rounded corners,minimum width=3.2cm,minimum height=1.05cm,align=center,fill=orange!10},
  flow/.style={-{Latex[length=2.5mm]},thick}
}
